\documentclass[10.5pt,a4paper,fontset=none]{ctexart}

\usepackage{geometry}
\usepackage{array}
\usepackage{booktabs}
\usepackage{enumitem}
\usepackage{setspace}
\usepackage{titlesec}
\usepackage{amsmath}
\usepackage{amssymb}
\usepackage{caption}
\usepackage{multirow}
\usepackage{xcolor}
\usepackage{fontspec}
\usepackage{float}
\usepackage[section]{placeins}
\usepackage{needspace}

\setCJKmainfont{Noto Serif CJK SC}
\setCJKsansfont{Noto Sans CJK SC}
\setCJKmonofont{Noto Sans Mono CJK SC}

\geometry{
  left=2.2cm,
  right=2.2cm,
  top=2.1cm,
  bottom=2.2cm
}

\pagestyle{empty}
\setlength{\parindent}{0pt}
\setlength{\parskip}{0.25em plus 0.04em minus 0.02em}
\linespread{1.25}
\setlength{\abovedisplayskip}{0.45em}
\setlength{\belowdisplayskip}{0.45em}
\setlength{\abovedisplayshortskip}{0.25em}
\setlength{\belowdisplayshortskip}{0.35em}
\setlength{\textfloatsep}{1.0em}
\setlength{\floatsep}{0.8em}
\setlength{\intextsep}{0.9em}
\captionsetup[table]{skip=5pt}
\renewcommand{\arraystretch}{1.08}
\setlength{\tabcolsep}{5pt}
\renewcommand{\topfraction}{0.95}
\renewcommand{\bottomfraction}{0.9}
\renewcommand{\textfraction}{0.05}
\renewcommand{\floatpagefraction}{0.85}

\titleformat{\section}
  {\bfseries\large}
  {\thesection}
  {1em}
  {}
\titleformat{\subsection}
  {\bfseries\zihao{4}}
  {\thesubsection}
  {0.8em}
  {}
\titlespacing*{\section}{0pt}{1.0em}{0.5em}
\titlespacing*{\subsection}{0pt}{0.8em}{0.38em}

\setlist[itemize]{
  leftmargin=2em,
  itemsep=0.2em,
  topsep=0.15em
}

\setlist[enumerate]{
  leftmargin=2.5em,
  itemsep=0.2em,
  topsep=0.15em
}

\begin{document}

\zihao{-4}

% ================= 封面信息 =================

\vspace*{0.8cm}

\begin{center}
  {\Large \textbf{多种算法求解0/1背包问题的性能比较}}
\end{center}

\vspace{0.5cm}

\begin{center}
\begin{tabular}{ccc}
\cline{1-3}
学号 & 姓名 & 任务分工 \\
\cline{1-3}
3024244487 & 张诗淇 & 回溯算法、分支限界算法、FSU \& Pi \& JJ 实验分析 \\
\cline{1-3}
3024244510 & 贾静潇 & 贪心算法、动态规划算法、数据整理与实验框架 \\
\cline{1-3}
\end{tabular}
\end{center}

\vspace{0.3cm}

\begin{center}
  \textbf{选题类型：}2.3 问题求解类 \quad \textbf{问题：}A.1 背包问题
\end{center}

\vspace{0.5cm}

% ================= 摘要 =================

\section{摘要}

本实验针对0/1背包问题，实现了贪心算法、动态规划算法、回溯算法和分支限界算法，并在FSU、Pi和JJ三个数据源上进行对比实验。FSU数据用于验证算法正确性，Pi生成数据用于观察不同规模和价值--重量相关性对算法表现的影响，JJ数据用于测试较大规模实例下的实际可运行性。实验结果表明：动态规划在所有可运行实例上均得到最优解，但随着容量和物品数量增加，运行时间明显上升；贪心算法速度最快，在多数大规模实例上接近最优，但在FSU P06中与最优值相差14.81\%；回溯和分支限界在小规模数据上能够求得最优解，但在$n=500$及以上规模中搜索空间过大，不适合继续完整运行。总体来看，四类算法在求解精度、时间开销和适用规模上各有侧重。

% ================= 实验目的 =================

\section{实验目的}

本实验希望通过0/1背包问题完成一次较完整的算法求解流程。具体目标包括：
\begin{enumerate}
  \item 理解0/1背包问题的数学模型和组合优化特征，掌握将实际问题抽象为算法模型的过程。
  \item 实现贪心、动态规划、回溯、分支限界四种算法，并保证同一数据格式下能够统一运行。
  \item 在不同来源、不同规模的数据集上比较算法的求解值、运行时间和稳定性。
  \item 从复杂度、最优性、数据分布和工程可运行性等角度归纳各算法的优缺点。
\end{enumerate}

% ================= 问题定义 =================

\section{问题定义}

给定$n$个物品，每个物品$i$有重量$w_i$和价值$v_i$，背包容量为$W$。每个物品只能选择一次，不能拆分。设$x_i \in \{0,1\}$表示是否选择第$i$个物品，则0/1背包问题可表示为：
\begin{equation}
\max \sum_{i=1}^{n} v_i x_i
\end{equation}
\begin{equation}
\text{s.t.} \quad \sum_{i=1}^{n} w_i x_i \leq W,\quad x_i \in \{0,1\}
\end{equation}

0/1背包问题是经典NP-hard问题。它既可以用动态规划在容量可控时求精确解，也可以用贪心、回溯、分支限界等策略从不同角度求解。由于这些算法的复杂度、剪枝方式和最优性保证差异明显，该问题适合作为本次多算法比较实验的对象。

从实验角度看，该问题有两个值得关注的参数：一是物品数量$n$，它直接影响组合搜索算法的搜索树规模；二是背包容量$W$，它直接影响动态规划的状态空间。价值与重量之间的相关性也会影响贪心和分支限界的表现。若单位价值差异明显，贪心排序往往更有效；若大量物品性价比接近，局部选择更容易与全局最优组合发生冲突。

% ================= 实验设计流程 =================

\section{实验设计流程与代码实现}

\subsection{总体流程}

本实验的整体流程分为五步。第一步，整理FSU、Pi、JJ三类数据，使程序能够统一读取物品数量、背包容量、物品重量和物品价值。第二步，实现四种算法，并尽量保持输入输出接口一致。第三步，对每个数据集运行实验入口程序，记录算法结果、官方最优值或基准最优值、运行时间和运行状态。第四步，将结果写入CSV总表，并同步生成按算法和数据集命名的单独结果文件，便于后续检查和制表。第五步，根据实验数据分析算法的复杂度特征、适用场景和局限性。

实验结果采用“总表+单独结果文件”的方式保存。总表为\texttt{experiments/result-summary.csv}，用于汇总全部实验记录，字段如下：
\[
\text{dataset\_source, dataset\_name, n, capacity, algorithm, value, optimum, gap\_percent, time\_ms, status}
\]
同时，程序会为每条实验记录生成单独的文本结果文件，命名格式为\texttt{算法-knapsack01-数据源-实例-c容量.txt}，例如\texttt{dp-knapsack01-fsu-p01-c165.txt}。单独结果文件采用键值对格式，便于检查某一次具体实验的参数和输出。根目录下的\texttt{readme.md}和\texttt{member.txt}也按照课程模板整理，分别说明代码、数据、结果文件和小组分工。

其中，\texttt{OK}表示算法正常完成，\texttt{SKIP}表示因规模过大主动跳过，\texttt{TIMEOUT}表示超过预设时间限制。这样设计的原因是：对于回溯和分支限界这类指数级搜索算法，不能简单要求它们在所有大规模数据上跑完，否则实验会陷入长时间无结果的状态。记录\texttt{SKIP}或\texttt{TIMEOUT}本身也是算法适用性分析的一部分。

\subsection{代码结构}

项目核心代码位于\texttt{src}目录。四个算法类分别为：
\begin{itemize}
  \item \texttt{GreedyKnapsack.java}：按价值重量比排序的贪心算法。
  \item \texttt{DPKnapsack.java}：一维数组优化的动态规划算法。
  \item \texttt{BacktrackingKnapsack.java}：带上界剪枝和超时保护的回溯算法。
  \item \texttt{BranchBoundKnapsack.java}：基于优先队列和上界剪枝的分支限界算法。
\end{itemize}

实验入口包括\texttt{Experiment.java}、\texttt{PiExperiment.java}和\texttt{JJExperiment.java}，分别用于运行FSU、Pi和JJ数据集。为了避免不同入口输出格式不一致，CSV汇总结果和单独文本结果均由\texttt{ExperimentCsvWriter.java}统一完成。Pi数据由\texttt{PiDataGenerator.java}生成，并固定随机种子，以保证不同成员和不同运行次数生成的数据一致。由于本项目使用Java语言，源文件名需要与类名保持一致，因此保留\texttt{GreedyKnapsack.java}、\texttt{DPKnapsack.java}等命名；这些文件名本身也体现了“算法名称+背包问题”的含义。

\subsection{数据集选择与实验策略}

FSU数据集包含8个标准小规模实例，适合检验算法结果是否能达到官方最优值。Pi数据集由程序生成，包含不相关、弱相关、强相关三种类型，每类设置S、M、L三种规模，共9个实例。JJ数据集从原始大规模数据中选取6个代表实例，覆盖$n=400$和$n=1000$两种规模，以及$f=0.1,0.2,0.3$三种参数。JJ完整数据较大，因此仓库中只保留下载脚本\texttt{setup\_jj.sh}，实际实验时下载指定实例，保证小组成员使用同一批数据。

\begin{table}[!htbp]
\centering
\footnotesize
\caption{实验数据集概况}
\begin{tabular}{ccccc}
\toprule
数据源 & 实例数 & 规模范围 & 容量范围 & 实验作用 \\
\midrule
FSU & 8 & $n=5$--24 & 26--6,404,180 & 验证算法正确性 \\
Pi & 9 & $n=50,500,2000$ & 约1.2万--254万 & 比较规模和相关性影响 \\
JJ & 6 & $n=400,1000$ & 1,000,000 & 测试大规模可运行性 \\
\bottomrule
\end{tabular}
\end{table}

实验策略上，FSU规模较小，因此四种算法全部运行。Pi数据中，S规模运行四种算法；M和L规模只运行贪心和DP，回溯与分支限界记录为\texttt{SKIP}。JJ数据中，物品数达到400或1000，回溯与分支限界也默认跳过。这样的安排不是减少实验内容，而是将复杂度理论反映到实验设计中：指数级搜索算法在大规模输入下不具备稳定可运行性，而贪心和DP可以继续比较解质量与运行时间。

% ================= 算法实现与复杂度 =================

\section{算法实现与复杂度分析}

\subsection{贪心算法}

贪心算法计算每个物品的单位重量价值$r_i=v_i/w_i$，按$r_i$从大到小排序，然后依次尝试装入背包。若当前物品重量不超过剩余容量，则加入；否则跳过该物品。这种策略来自分数背包问题，在物品可以拆分时能够得到最优解。但0/1背包中物品不可拆分，因此该策略只能作为启发式算法。

贪心算法的主要时间消耗来自排序，时间复杂度为$O(n\log n)$，空间复杂度为$O(n)$。它不需要枚举容量状态，也不需要遍历组合搜索树，因此对容量$W$基本不敏感。无论$W$是几百还是几百万，算法流程都主要由物品数$n$决定。

该算法的优点是实现简单、速度快，尤其适合大规模近似求解。当物品单位价值差异明显时，贪心排序通常能选出较好的解。但它只根据当前性价比做局部选择，不能回退或重新组合，所以不保证全局最优。本实验中，贪心在JJ数据上gap很小，但在FSU P06上gap达到14.81\%，说明它的效果与数据分布和物品组合结构密切相关。

\subsection{动态规划算法}

动态规划利用0/1背包的最优子结构。定义$dp[c]$为容量为$c$时能够获得的最大价值。对每个物品$i$，按照容量从大到小进行状态转移：
\begin{equation}
dp[c] = \max(dp[c], dp[c-w_i]+v_i), \quad c \geq w_i
\end{equation}
倒序遍历容量可以避免同一物品在一次循环中被重复使用。如果容量正序更新，当前物品可能被多次计入，问题会变成完全背包，这与0/1背包的约束不一致。

动态规划的时间复杂度为$O(nW)$，使用一维数组后空间复杂度为$O(W)$。这里的$W$是背包容量的数值，因此该算法属于伪多项式算法。它的复杂度并不是只与输入规模$n$有关，而是直接取决于容量大小。例如，当$n=2000$、$W$约为250万时，需要进行数十亿级别的状态更新，运行时间会明显上升。

动态规划的主要优点是结果稳定，能够保证最优解，适合容量适中的精确求解场景。本实验中DP在FSU、Pi和JJ全部可运行实例上均达到最优值。它的局限也很明显：复杂度直接依赖容量$W$，因此当$W$较大时，即使物品数量不多，也会导致时间或内存压力。例如FSU P08只有24个物品，但容量为6,404,180，DP耗时达到757.697ms，明显高于其他FSU实例。

\subsection{回溯算法}

回溯算法把每个物品看成一个二叉选择：选或不选。它从第一个物品开始递归搜索所有可行组合，并记录当前发现的最大价值。与DP不同，回溯不按容量建立状态表，而是直接在解空间树上搜索。每条从根到叶子的路径对应一种物品选择方案。

为了减少无效搜索，本实验在回溯中加入了两类剪枝。第一类是可行性剪枝：如果当前重量已经超过背包容量，则该分支不可能产生可行解，直接返回。第二类是上界剪枝：根据剩余物品估算该分支理论上最多还能获得多少价值，如果当前价值加上上界仍不超过已知最优值，就停止扩展该分支。上界估计越紧，剪枝效果越好。

不考虑剪枝时，回溯算法最坏需要枚举$2^n$种选择，时间复杂度为$O(2^n)$；递归深度最多为$n$，若只保存当前路径，空间复杂度为$O(n)$。回溯法的优点是逻辑直观，并且能够得到最优解，适合小规模实例或用于验证其他算法正确性。但它对物品数量$n$非常敏感，当$n$达到几百时，搜索空间已经远超合理运行范围。因此在Pi-M、Pi-L和JJ实例中，实验中将其记录为\texttt{SKIP}。

\subsection{分支限界算法}

分支限界算法同样基于搜索树，但它不一定按照深度优先顺序搜索，而是使用优先队列保存待扩展节点，并优先扩展上界更高的节点。每个节点包含当前层数、当前重量、当前价值和上界估计。当某个节点的上界不超过当前最优值时，该节点不再继续扩展。本实验使用类似分数背包的思想计算上界，使剪枝更有依据。

与回溯相比，分支限界更强调“有希望的节点优先”。这种策略通常可以更早找到较好的可行解，从而提高后续剪枝的效果。特别是在上界估计较准确时，大量不可能超过当前最优值的节点会被提前排除。

分支限界算法在剪枝有效时通常比普通回溯更快，但最坏情况下仍可能访问指数级节点，因此时间复杂度仍为$O(2^n)$。同时，由于优先队列可能保存大量活节点，最坏空间复杂度也可能达到$O(2^n)$。它的优点是能保证最优，且比单纯深度优先搜索更有方向性；缺点是实现复杂度较高，对上界估计质量和数据分布较敏感，在大规模数据上仍不具备稳定可运行性。

% ================= 实验结果 =================

\FloatBarrier
\Needspace{10\baselineskip}
\section{实验结果与分析}

\Needspace{6\baselineskip}
\subsection{FSU数据集结果}

\begin{table}[!htbp]
\centering
\footnotesize
\caption{FSU：四种算法结果对比}
\begin{tabular}{cccccccc}
\toprule
实例 & $n$ & 贪心值 & DP值 & 回溯值 & 分支限界值 & 最优值 & 贪心Gap(\%) \\
\midrule
P01 & 10 & 309 & 309 & 309 & 309 & 309 & 0.00 \\
P02 & 5 & 47 & 51 & 51 & 51 & 51 & 7.84 \\
P03 & 6 & 146 & 150 & 150 & 150 & 150 & 2.67 \\
P04 & 7 & 102 & 107 & 107 & 107 & 107 & 4.67 \\
P05 & 8 & 858 & 900 & 900 & 900 & 900 & 4.67 \\
P06 & 7 & 1478 & 1735 & 1735 & 1735 & 1735 & 14.81 \\
P07 & 15 & 1441 & 1458 & 1458 & 1458 & 1458 & 1.17 \\
P08 & 24 & 13415886 & 13549094 & 13549094 & 13549094 & 13549094 & 0.98 \\
\bottomrule
\end{tabular}
\end{table}

\begin{table}[!htbp]
\centering
\footnotesize
\caption{FSU：四种算法耗时对比（ms）}
\begin{tabular}{ccccc}
\toprule
实例 & 贪心 & DP & 回溯 & 分支限界 \\
\midrule
P01 & 23.052 & 14.113 & 44.950 & 30.434 \\
P02 & 0.066 & 0.029 & 0.067 & 0.091 \\
P03 & 0.026 & 0.065 & 0.035 & 0.072 \\
P04 & 0.032 & 0.038 & 0.037 & 0.066 \\
P05 & 0.028 & 0.057 & 0.200 & 0.088 \\
P06 & 0.161 & 0.069 & 0.070 & 0.271 \\
P07 & 0.084 & 0.921 & 0.573 & 0.760 \\
P08 & 0.081 & 757.697 & 0.306 & 1.378 \\
\bottomrule
\end{tabular}
\end{table}

FSU实例规模较小，因此四种算法都能完成运行。DP、回溯和分支限界在8个实例上均得到最优值，说明三种精确算法实现正确。贪心算法只有P01达到最优，其余实例均存在偏差，其中P06的gap达到14.81\%，是本实验中贪心误差最大的案例。这说明单位价值最高的物品不一定能组成整体最优方案，局部选择与全局组合之间可能存在冲突。

从运行时间看，P08最能体现容量对DP的影响。P08的物品数只有24，但容量超过640万，DP耗时757.697ms，远高于其他小容量实例。相比之下，回溯和分支限界主要受物品数和剪枝效果影响，因此在P08中仍能较快完成。这一组结果说明，算法复杂度中的主导参数并不相同：DP主要受$W$影响，而搜索类算法更受$n$影响。对于P01等极小实例，毫秒级结果容易受到JVM类加载和系统调度影响，不宜单独作为速度判断依据。

FSU结果还体现了“最优性”和“速度”不能简单等同。贪心在P06只用0.161ms，但结果比最优值少257；DP、回溯、分支限界耗时略高，但都能保证最优。因此，如果应用场景不能接受价值损失，应优先选择精确算法；如果只需要快速给出一个可行方案，贪心更有优势。

\FloatBarrier

\Needspace{8\baselineskip}
\subsection{Pi生成数据集结果}

\begin{table}[!htbp]
\centering
\footnotesize
\caption{Pi：四种算法结果对比}
\begin{tabular}{ccccccc}
\toprule
实例 & $n$ & 贪心值 & DP值 & 回溯值 & 分支限界值 & 贪心Gap(\%) \\
\midrule
UNCORR\_S & 50 & 19976 & 19976 & 19976 & 19976 & 0.00 \\
UNCORR\_M & 500 & 1001047 & 1001186 & SKIP & SKIP & 0.01 \\
UNCORR\_L & 2000 & 4111121 & 4111268 & SKIP & SKIP & 0.00 \\
WEAK\_S & 50 & 13822 & 13822 & 13822 & 13822 & 0.00 \\
WEAK\_M & 500 & 703546 & 703810 & SKIP & SKIP & 0.04 \\
WEAK\_L & 2000 & 2866292 & 2866401 & SKIP & SKIP & 0.00 \\
STRONG\_S & 50 & 16297 & 16437 & 16437 & 16437 & 0.85 \\
STRONG\_M & 500 & 800642 & 801635 & SKIP & SKIP & 0.12 \\
STRONG\_L & 2000 & 3199654 & 3200355 & SKIP & SKIP & 0.02 \\
\bottomrule
\end{tabular}
\end{table}

\begin{table}[!htbp]
\centering
\footnotesize
\caption{Pi：四种算法耗时对比（ms）}
\begin{tabular}{ccccc}
\toprule
实例 & 贪心 & DP & 回溯 & 分支限界 \\
\midrule
UNCORR\_S & 23.471 & 16.967 & 23.391 & 34.255 \\
UNCORR\_M & 1.299 & 528.519 & SKIP & SKIP \\
UNCORR\_L & 4.291 & 10056.173 & SKIP & SKIP \\
WEAK\_S & 0.173 & 3.882 & 0.689 & 1.092 \\
WEAK\_M & 0.939 & 657.953 & SKIP & SKIP \\
WEAK\_L & 3.115 & 10663.162 & SKIP & SKIP \\
STRONG\_S & 0.083 & 1.558 & 0.218 & 1.549 \\
STRONG\_M & 4.324 & 739.220 & SKIP & SKIP \\
STRONG\_L & 4.461 & 11416.532 & SKIP & SKIP \\
\bottomrule
\end{tabular}
\end{table}

Pi数据更适合观察规模变化。S规模$n=50$时，四种算法都能运行，DP、回溯和分支限界均达到最优。贪心在不相关和弱相关S实例上也达到最优，但在强相关S实例上出现0.85\%的gap。原因在于强相关数据中物品价值与重量变化趋势接近，单位价值差异较小，排序策略的区分度下降。此时前几个被贪心选中的物品未必能与剩余容量形成最优组合。

随着规模增大到M和L，回溯与分支限界被记录为\texttt{SKIP}。这并不是程序缺失，而是由复杂度决定的实验处理方式。以$n=500$为例，二叉搜索空间已达到$2^{500}$量级，继续运行很难在合理时间内得到结果。DP仍能提供最优值，但耗时明显上升：不相关数据从S到L分别为16.967ms、528.519ms和10056.173ms。贪心则始终保持毫秒级运行，在Pi全部实例中最大gap为0.85\%，说明在该类生成数据上具有较好的近似效果。

从数据相关性看，贪心gap并不是简单随$n$增加而单调增大，而是与物品价值和重量的关系有关。强相关数据中，价值和重量同步变化，单位价值更接近，贪心排序的优势被削弱；不相关数据中，单位价值差异更大，贪心更容易识别高性价比物品。这说明评价贪心算法不能只看输入规模，还要关注数据分布。

从规模扩展看，DP的运行时间随$nW$增长明显上升。Pi-L三个实例的DP时间均超过10秒左右，而贪心仅需数毫秒。这一对比说明，如果目标是严格最优，DP是较可靠的选择；如果目标是快速得到接近最优的方案，贪心在该类数据上更符合工程需求。

\FloatBarrier

\Needspace{8\baselineskip}
\subsection{JJ大规模数据集结果}

\begin{table}[!htbp]
\centering
\footnotesize
\caption{JJ：四种算法结果对比}
\begin{tabular}{ccccccc}
\toprule
实例 & $n$ & 贪心值 & DP值 & 回溯值 & 分支限界值 & 贪心Gap(\%) \\
\midrule
n=400,f=0.1 & 400 & 1003879 & 1003939 & SKIP & SKIP & 0.01 \\
n=400,f=0.2 & 400 & 1004227 & 1004245 & SKIP & SKIP & 0.00 \\
n=400,f=0.3 & 400 & 1004587 & 1005007 & SKIP & SKIP & 0.04 \\
n=1000,f=0.1 & 1000 & 1007973 & 1008317 & SKIP & SKIP & 0.03 \\
n=1000,f=0.2 & 1000 & 1008470 & 1008509 & SKIP & SKIP & 0.00 \\
n=1000,f=0.3 & 1000 & 1010292 & 1010297 & SKIP & SKIP & 0.00 \\
\bottomrule
\end{tabular}
\end{table}

\begin{table}[!htbp]
\centering
\footnotesize
\caption{JJ：四种算法耗时对比（ms）}
\begin{tabular}{ccccc}
\toprule
实例 & 贪心 & DP & 回溯 & 分支限界 \\
\midrule
n=400,f=0.1 & 27.783 & 788.311 & SKIP & SKIP \\
n=400,f=0.2 & 0.516 & 750.138 & SKIP & SKIP \\
n=400,f=0.3 & 0.990 & 727.072 & SKIP & SKIP \\
n=1000,f=0.1 & 2.051 & 1556.003 & SKIP & SKIP \\
n=1000,f=0.2 & 1.860 & 1616.109 & SKIP & SKIP \\
n=1000,f=0.3 & 0.959 & 1632.430 & SKIP & SKIP \\
\bottomrule
\end{tabular}
\end{table}

JJ实验展示了较大规模下算法的实际差异。DP在6个实例上均得到最优值，说明在容量为1,000,000且$n \leq 1000$时，一维DP仍然可以运行。但耗时随物品数上升明显增加：$n=400$时约727--788ms，$n=1000$时约1556--1632ms。贪心算法在JJ上表现很好，gap最大仅0.04\%，且除首次运行实例外基本保持在毫秒级。

JJ结果说明，在某些结构化大规模数据中，贪心虽然没有理论最优保证，但实际结果可能非常接近最优。其原因可能是这些实例中物品价值和重量的关系使单位价值排序具有较好的指导作用。不过，这一现象不能推广为贪心总是可靠，因为FSU P06已经给出了反例。较稳妥的结论是：贪心适合先给出快速基线，但如果需要严格最优，仍需要DP或其他精确算法验证。

回溯和分支限界在JJ上均跳过。对于$n=400$或$n=1000$的实例，即使剪枝有效，也无法保证在合理时间内完成指数级搜索。这个结果说明，精确搜索算法虽然理论上能得到最优解，但可运行性受到规模限制；在大规模场景中，贪心和动态规划更具有现实意义。

\FloatBarrier

% ================= 综合分析 =================

\section{综合分析}

\Needspace{8\baselineskip}
\subsection{复杂度与实验现象对照}

\begin{table}[!htbp]
\centering
\footnotesize
\caption{四种算法复杂度与适用性对比}
\begin{tabular}{p{2.0cm}p{2.6cm}p{2.7cm}p{2.7cm}p{2.8cm}}
\toprule
维度 & 贪心 & 动态规划 & 回溯 & 分支限界 \\
\midrule
时间复杂度 & $O(n\log n)$ & $O(nW)$ & 最坏$O(2^n)$ & 最坏$O(2^n)$ \\
空间复杂度 & $O(n)$ & $O(W)$ & $O(n)$ & 最坏$O(2^n)$ \\
是否保证最优 & 否 & 是 & 是 & 是 \\
主要敏感因素 & 数据分布 & 容量$W$ & 物品数$n$ & 物品数$n$、上界紧度 \\
主要优点 & 快速、实现简单 & 稳定、可作最优基准 & 思路直观、能求最优 & 剪枝更强、搜索更有方向 \\
主要缺点 & 不保证最优 & 大容量下开销大 & 大规模搜索爆炸 & 仍可能指数爆炸，内存压力较大 \\
适用场景 & 大规模近似求解 & 容量适中的精确求解 & 小规模精确求解 & 小到中等规模且剪枝有效的精确求解 \\
\bottomrule
\end{tabular}
\end{table}

从实验结果看，理论复杂度和实际表现基本一致。贪心算法在Pi-L和JJ中仍能快速完成，说明$O(n\log n)$的可扩展性较好；DP在所有实例上结果最稳定，但P08和Pi-L中的运行时间明显上升，体现出$O(nW)$对容量的依赖；回溯和分支限界在FSU与Pi-S中表现正常，但在$n=500$及以上实例中被跳过，反映出指数复杂度的实际限制。

\Needspace{8\baselineskip}
\subsection{最优性与运行时间的折中}

四种算法可以看作位于一条“速度--精度”折中线上。贪心位于速度端，运行极快但不保证最优；DP位于稳定精确端，只要容量可接受就能给出最优；回溯和分支限界也能保证最优，但更依赖物品数量和剪枝效果。实验结果表明，单纯比较时间或单纯比较最优性都不完整，算法选择需要结合应用要求。

如果一个场景要求严格最优，例如预算分配、资源受限且价值损失不可接受，则DP更合适。如果场景更关注响应速度，例如需要快速生成候选方案或实时估计，贪心更实用。回溯和分支限界适合小规模精确求解，也适合在教学实验中观察搜索树和剪枝的作用，但不适合直接处理几百或几千个物品的实例。

\Needspace{8\baselineskip}
\subsection{容量、物品数和数据分布的影响}

本实验还说明，不同算法受不同因素制约。DP主要受容量$W$制约，FSU P08就是典型例子：$n$只有24，但由于$W$很大，DP耗时显著增加。回溯和分支限界主要受物品数$n$制约，Pi-M和JJ虽然容量不一定比P08更极端，但$n$达到数百或上千，搜索空间已不可接受。贪心主要受数据分布影响，FSU P06的较大gap说明某些组合结构会让局部性价比选择失效。

因此，不能简单说某一种算法“最好”。同一个0/1背包问题，在小规模、容量适中、要求最优的情况下适合DP或搜索；在大规模、允许近似的情况下适合贪心；在需要研究搜索过程和剪枝效果时适合回溯或分支限界。算法设计的关键不是记住某个固定答案，而是根据输入特征选择合适策略。

\Needspace{8\baselineskip}
\subsection{实验局限与改进方向}

本实验仍有一些局限。首先，运行时间受Java虚拟机预热和系统状态影响，极小实例上的毫秒级差异不一定完全稳定。其次，Pi数据虽然固定了随机种子，但仍属于生成数据，不能覆盖所有真实场景。再次，回溯和分支限界在大规模实例中采用\texttt{SKIP}，虽然符合复杂度分析，但没有进一步比较不同剪枝策略的差异。

后续可以从几个方面改进：一是多次重复运行取平均时间，减少偶然波动；二是为贪心增加不同排序策略，例如按价值降序、按重量升序，与单位价值策略比较；三是进一步优化分支限界的上界估计和初始可行解；四是尝试近似算法或启发式算法，如局部搜索、遗传算法等，用于补充DP和贪心之间的性能空档。

\FloatBarrier

% ================= 总结 =================

\section{总结}

本实验完成了0/1背包问题的多算法求解与性能比较，实现了贪心、动态规划、回溯和分支限界四种方法，并在FSU、Pi、JJ三类数据集上进行了测试。实验表明，动态规划结果最稳定，但受容量影响明显；贪心速度最快，但可能偏离最优；回溯和分支限界在小规模数据上能得到最优解，但面对大规模实例时搜索空间过大。通过本次实验，我们更直观地理解了算法复杂度对实际运行效果的影响，也认识到算法选型需要综合考虑数据规模、容量大小、解的精度要求和运行时间限制。

\end{document}
